\scanpage{13}
\mspage{\textcircled{3}}
% Reversed typescript shows through.
\uline{NB} Le rev. universel de $G$ est de \deleted{\unsure{groupe}}\inserted{noyau} $\mathbb Z$. Il donne une extension \inserted{centrale} de $\Gamma$ par $\mathbb Z$, d'où un élément de $H^2(\Gamma,\mathbb Z)\simeq H^2(X,\mathbb Z)$. Montrons que c'est $\gamma_1$. Plus généralement, \unsure{quand on a|partant on a} une extension d'un gpe top. $G$ par un groupe discret $H$, lui correspond un elt de $H^2(B_G,H)$ (\unsure{fonctorielt|fonctoriel} en $G,H$ \ldots).
\margin{\diagram{\begin{tikzcd}[ampersand replacement=\&,column sep=tiny]
B_H\arrow[r]\&B_E\arrow[r]\&B_G
\end{tikzcd}}}
Dans le cas où $G$ est $S=$ tore ``\unsure{standart|standard}, et $\widetilde S=$ rev. universel de $\widetilde S$, de noyau $\mathbb Z$, l'\unsure{élém.} qu'on \unsure{trouve} \unsure{dans} $H^{\unsure{2|*}}(B_S,\mathbb Z)$ n'est autre que le générateur canonique de Chern.

Ces remarques ensemble \unsure{établissent|établ\textsuperscript{t}} \unsure{expmt} l'interprétation \unsure{signalée|esquissée} de $\gamma_1^X\in$ \deleted{\unsure{dans}} $H^2(X,\mathbb Z)\simeq H^2(\Gamma,\mathbb Z)$. Donc son image dans $H^2(\Gamma,\mathbb Z/2\mathbb Z)$ \deleted{\unsure{est nulle}}\inserted{$=\mathbb Z/2\mathbb Z$} est la classe de l'extension de $\Gamma$ par $\mathbb Z/2\mathbb Z$, déduite de $\widetilde G\longrightarrow G$. \unsure{C'est nulle|C'est celle} donc \deleted{\unsure{ssi}} \unsure{précisément} \unsure{si scindée|si elle}. Comme $\mathrm{EP}(X)=2(1-g)$ est pair, elle est tjrs nulle !

On gagne ! Il existe tjrs un $\widetilde\Gamma$ !

\textcircled{4} Revenant, avec cette information, à la situation \textcircled{1}, on voit que \deleted{\unsure{la}}, pour tte courbe alg. \unsure{(non complète|(non complétée|lisse complète} connexe de genre $g\geq2$, fibré tangent $T_X$, le fibré vectoriel \deleted{complexe}\inserted{réel de rang $\unsure{4|2|1}$} $T_X\deleted{\unsure{\otimes_{\mathbb R}T_X}}+2\mathbb{1}_{\mathbb R}$ sur $X$ est muni d'une connexion intégrable, (attention, elle ne respecte pas les
% archive sheet number: 12
